% AUTORIDAD DE FUENTE: Exposé IX §12.8 hasta antes de (12.8.9), líneas
% francesas 3709--3795; índices de escaneo 488--491, folios 477--480.

\medskip
\noindent
12.8. \underline{Demostración de 12.5: reducción al caso en que $S$ es
de característica cero.} Utilizaremos un razonamiento basado en el mismo
principio que la demostración del «teorema sobre la acción esencialmente
moderada de la monodromía sobre $\pi_1$» en el Exposé
V\sisi{\hspace{4em}}{}. Obsérvese que, por 10.3, podemos suponer que $S$
es completo. Sea $W$ un $p$-anillo de Cohen con cuerpo residual $k$, y
elijamos un homomorfismo $W\longrightarrow V$ que induzca la identidad
en los cuerpos residuales (EGA $0_{\mathrm{IV}}$, 19.8.6). Sea
$(x_\lambda)_{\lambda\in L}$ el conjunto de los puntos singulares de
$X_o$, y consideremos el esquema formal de módulos $\underline S$ de
$X_o$ sobre $W$, junto con la «curva universal»
\[
 \underline f:\underline X\longrightarrow\underline S
\]
cuya fibra especial es $X_o$ (Exposé VI\sisi{\hspace{3em}}{}). Se vio en
el lugar citado que $\underline S$ es isomorfo al espectro formal de un
anillo de series formales $W[[T_1,\ldots,T_N]]$, y que en $\underline S$
puede hallarse un divisor con cruces normales

%%%%%%%%%%  scan idx 489  |  folio 478  %%%%%%%%%%

\begin{equation}\tag{12.8.1}
 D=\sum_{\lambda\in L}D_\lambda,
\end{equation}

\noindent
cuyas componentes irreducibles $D_\lambda$ están indexadas por $L$, y
cada $D_\lambda$ está descrita por una ecuación de la forma
$T_{i_\lambda}=0$. (Si se desea, puede suponerse que
$i_\lambda=\lambda$ tras indexar $L$ por un conjunto de enteros
$[1,r]$.) La componente $D_\lambda$ es la imagen por
$\underline f:\underline X\longrightarrow\underline S$ de la componente
$D'_\lambda$ que pasa por el punto singular $x_\lambda$ de $X_o$ en el
lugar de no lisura
\[
 D'=\bigcup_{\lambda\in L}D'_\lambda
\]
de $\underline f$. Por la propia construcción,
$f:X\longrightarrow S$ es isomorfo a la imagen inversa de
$\underline f:\underline X\longrightarrow\underline S$ por un
$W$-morfismo adecuado $g:S\longrightarrow\underline S$:

\begin{equation}\tag{12.8.2}
\begin{tikzcd}
 \underline X \arrow[d,"\underline f"'] &
 X \arrow[l] \arrow[d,"f"] \\
 \underline S & S \arrow[l,"g"']
\end{tikzcd}
\qquad .
\end{equation}

\noindent
Pongamos

\begin{equation}\tag{12.8.3}
 P=\underline{\operatorname{Pic}}^{o}_{\underline X/\underline S}.
\end{equation}

\noindent
Como la fibra especial de $\underline f$ es geométricamente reducida de
dimensión $1$, se sigue que $\sisi{\underline F}{\underline f}$ es cohomológicamente plano en
dimensión cero, por EGA III, 7.8.1, y que $P$ es formalmente liso sobre
$\underline S$ (esto último resulta de
$H^2(X_o,\underline{\mathcal O}_{X_o})=0$, puesto que $X_o$ tiene
dimensión $1$). Para simplificar, supondremos que $P$
\underline{es representable por un esquema en grupos sobre
$\underline S$}; será necesariamente un esquema en grupos liso de
fibras conexas. Este es un caso particular de un resultado no publicado
de Mumford citado en [22] (y considerablemente generalizado por Raynaud
en la nota allí citada, cuando la base es un trait). Como no se ha
publicado ninguna demostración de los resultados de Mumford o de Raynaud,
el lector puede preferir utilizar los resultados generales de M. Artin
[1]. Estos proporcionan la representabilidad de $P$---bajo la única
hipótesis de que $\underline f$ sea propio, plano y cohomológicamente
plano en dimensión cero---por un «espacio algebraico» en el sentido de
Artin, lo cual bastaría para nuestro propósito. Pongamos

%%%%%%%%%%  scan idx 490  |  folio 479  %%%%%%%%%%

\begin{equation}\tag{12.8.4}
 U=\underline S-\operatorname{supp}D.
\end{equation}

\noindent
Así, $U$ es el mayor abierto de $\underline S$ sobre el cual
$\underline X$ es liso, y, por consiguiente,
\[
 P|_U=\underline{\operatorname{Pic}}^{o}_{\underline X_U/U}
\]
es un esquema abeliano sobre $U$. Su imagen inversa sobre $\eta$ por
$g:\eta\longrightarrow U$ es el esquema abeliano $A_\eta$ sobre $\eta$
que queremos estudiar. Más precisamente, existe un isomorfismo canónico
sobre $S$ (12.1.10):

\begin{equation}\tag{12.8.5}
 P\times_{\underline S}S\simeq A^{o}.
\end{equation}

\noindent
Como la fibra especial $P_o=A_o^{o}$ es $\ell$-divisible para todo
$\ell$, lo mismo ocurre con $P$; todas sus fibras son, por tanto,
extensiones de una variedad abeliana por un toro. Consideremos el grupo
pro-$\ell$ $T_\ell(P)$ sobre $\underline S$, cuya imagen inversa por $g$
es precisamente $T_\ell(A^{o})$. Aplicando las construcciones de 7.3, se
define una filtración de este progrupo mediante una «parte fija» y una
«parte tórica»:

\begin{equation}\tag{12.8.6}
 T_\ell(P)\supset T_\ell(P)^f=T_\ell(\widehat P)
 \supset T_\ell(P)^t=T_\ell(T)\supset0.
\end{equation}

\noindent
Su imagen inversa por $g$ es la filtración análoga de $T_\ell(A^{o})$
estudiada en las secciones 2 y 5.

%%%%%%%%%%  scan idx 491  |  folio 480  %%%%%%%%%%

Consideremos la restricción de esta filtración al abierto $U$ de
$\underline S$. Allí tenemos la autodualidad canónica de $P_U$, por ser
la jacobiana de un esquema de curvas liso sobre $U$, y, por tanto, el
emparejamiento correspondiente

\begin{equation}\tag{12.8.7}
 \varphi:T_\ell(P_U)\otimes T_\ell(P_U)
 \longrightarrow T_\ell(\mathbb G_{mU}).
\end{equation}

\noindent
Esta es una \underline{autodualidad}, en el sentido de los grupos de
Barsotti--Tate, de $T_\ell(\sisi{B_U}{P_U})$. Se vio en 7.4.8 que, para
este emparejamiento, $T_\ell(P_U)^f$ \underline{y} $T_\ell(P_U)^t$
\underline{son ortogonales uno del otro}. Se sigue que
\[
 E_U=T_\ell(P_U)/T_\ell(P_U)^t
\]
es canónicamente isomorfo al progrupo de Barsotti--Tate dual de
$T_\ell(P_U)^f=F_U$. Se prolonga, pues, canónicamente a un grupo de
Barsotti--Tate $E$ sobre $\underline S$, a saber,

\begin{equation}\tag{12.8.8}
 E=D(F),\qquad\text{donde}\qquad F=T_\ell(P)^f.
\end{equation}

\noindent
Nos encontramos ahora en una situación típica de extensiones mixtas
(9.3), y podemos repetir las construcciones de 9.4--9.6 (¡por haberlas
escrito en el marco estrecho y de mierda de un esquema base reducido a
un trait, pedimos disculpas!): $T_\ell(P_U)$ es una extensión mixta de
$E_U$ por $F_U$, donde $E$ y $F$ son grupos de Barsotti--Tate sobre
$\underline S$ que son extensiones duales entre sí:
